\begin{textsample}{POS Dim 1 – vem\_ed\_07 – Score 11.00 – vem\_ed\_07/t027.txt}  \label{ex:f1_pos_183}
Vozes femininas da literatura brasileira O Colóquio de Literatura Brasileira \textbf{foi} um evento organizado pelas Fatecs Guaratinguetá e Itaquaquecetuba com a Uniminuto, em 14 de maio de 2021, e transmitido pela plataforma Lark.
Voltado para estudantes de Licenciatura em Línguas Estrangeiras da universidade colombiana, o colóquio \textbf{foi} moderado por Osvaldo Succi Junior, coordenador dos PCIs / Cesu. \textbf{Contou} com a participação dos diretores da Fatec Guaratinguetá, André Amarante, e da Fatec Itaquaquecetuba, Sonia Maria Alvarez.
Desta unidade, marcaram presença ainda os coordenadores Francisco Claudio Tavares, do curso de Gestão Comercial, e Aparecido Rodrigues da Silva López-Guerrero, de Gestão da Tecnologia da Informação.
Regiane Souza Camargo Moreira (Fatec Guaratinguetá) palestrou sobre"O feminino como representação na Literatura: do Romantismo ao Modernismo".
Wilton Garcia (Fatec Itaquaquecetuba) apresentou"Literatura Contemporânea: vozes negras, vozes femininas".
Destacou autoras como Sueli Carneiro, fundadora do Geledés - Instituto da Mulher Negra; Djamila Ribeiro, filósofa e feminista negra; e Carolina Maria de Jesus, autora de Quarto de Despejo.
Alertou sobre o silenciamento das vozes das mulheres negras na literatura e na sociedade.
Javier Guerrero, coordenador de \textbf{desenvolvimento} de professores da Uniminuto, acompanhou o Colóquio.
Sobre as parcerias com as Fatecs, comenta:"entre Brasil e Colômbia temos muitas \textbf{coisas} em \textbf{comum}, mas também \textbf{diferenças}, como a língua, e contrastes \textbf{que} nos \textbf{permitem} \textbf{aprender}, e isso \textbf{é} perfeito para as colaborações".
QUEM \textbf{É} QUEM James Goodwin \textbf{é} diretor de \textbf{desenvolvimento} estudantil e iniciativas acadêmicas da BridgeValley Community and Technical College (EUA).
No segundo semestre de 2019, a instituição \textbf{iniciou} um projeto \textbf{colaborativo} (COIL) sobre equipamentos médico-hospitalares com a Fatec Bauru (leia mais na p.4).
Desde 2020, desenvolve um projeto sobre física das ondas, com estudantes de Produção Fonográfica da Fatec Tatuí.
Atualmente participam 18 alunos dos EUA e 27 brasileiros, orientados pelos professores Pedro Rosa e Dulce Villa Nova (Fatec Tatuí) e Machele Kindle (BridgeValley).
Outros dois projetos estão em andamento em Bridge Valley: com o Symbiosis Center for Management Studies (Índia), com 14 alunos norte-americanos e 28 indianos; e com o La Guardia Community College (EUA), com 12 estudantes de BridgeValley e 20 de La Guardia.
Um professor de \textbf{cada} instituição orienta esses COILs."A nossa história com os COILs \textbf{começou} em 2019, com a consultoria de Jon Rubin", \textbf{conta} Goodwin.
Rubin \textbf{foi} o criador da iniciativa COIL na State University of New York, em 2006.
Para os próximos anos, os objetivos da BridgeValley \textbf{são}:"primeiro, aumentar o \textbf{número} de professores oferecendo COILs e, consequentemente, de alunos participando; e segundo, integrar os COILs ao currículo, como \textbf{parte} dos requisitos de educação \textbf{geral} nos cursos de graduação".
Goodwin espera ainda \textbf{ampliar} as parcerias com as Fatecs,"não só nos modelos tradicionais de projetos entre faculdades, mas também com conversas culturais, discussões e \textbf{intercâmbios} virtuais administrativos".

% matched lemmas: ampliar, aprender, cada, coisa, colaborativo, começar, comum, contar, desenvolvimento, diferença, geral, iniciar, intercâmbio, número, parte, permitir, que, ser
\end{textsample}
